\providecommand{\Z}{\mathbb Z}
\providecommand{\C}{\mathbb C}
\providecommand{\La}{\Lambda}
\providecommand{\Perf}{\operatorname{Perf}}
\providecommand{\RPerf}{\operatorname{RPerf}}
\providecommand{\Coh}{\operatorname{Coh}}
\providecommand{\RHom}{\operatorname{RHom}}
\providecommand{\Ext}{\operatorname{Ext}}
\providecommand{\sdet}{\operatorname{sdet}}
\providecommand{\Supp}{\operatorname{Supp}}
\providecommand{\Ran}{\operatorname{Ran}}
\providecommand{\CoHA}{\operatorname{CoHA}}
\providecommand{\Hall}{\mathcal H}
\providecommand{\Prim}{\operatorname{Prim}}
\providecommand{\sdim}{\operatorname{sdim}}
\providecommand{\Spch}{\mathrm{Sp}^{\mathrm{ch}}}
\providecommand{\hCS}{\mathrm{hCS}}
\providecommand{\BV}{\mathrm{BV}}
\providecommand{\TS}{\mathrm{TS}}
\providecommand{\Mot}{\mathrm{Mot}}
\providecommand{\epsor}{\epsilon_o}
\providecommand{\nuDelta}{\nu_{\Delta_5}}
\providecommand{\gDelta}{\mathfrak g_{\Delta_5}}

\section{Agent 15-C: oriented compact Hall/CoHA atlas attack--heal}
\label{sec:agent15c-oriented-compact-hall-coha-atlas}

\subsection{Attacked claim}

The attacked claim is the strong gate-3 assertion:
\[
  \hbox{the compact oriented Hall/CoHA atlas for }K3\times E
  \hbox{ already exists in the manuscript or in the adjacent volumes.}
\]
This is false.  The current Igusa manuscript, the six Agent 14 reports,
and the surrounding volumes contain the correct typed requirements,
finite-stage comparison technology, and several scalar or quotient
shadows.  They do not contain the actual extension-closed compact
moduli stacks, mixed wrapped extension correspondences, reduced
\(E\)-quotient descent, strong orientation system, protected
integration map, Weyl lift, and primitive presentation comparison.

The healed statement is therefore:

\begin{quote}
The oriented compact Hall/CoHA atlas is a conditional construction
problem.  If the exact moduli objects, critical charts, orientations,
finite-first HN truncations, mixed local/wrapped correspondences, and
proper pull--push formalism listed below are supplied, then they assemble
to the gate-3 pro-Hall object required by the compact Igusa realization
datum.  None of these data is constructed by the scalar identity
\(-4096\Delta_5^{-2}\), by the Borcherds product, or by target-internal
bar--cobar inversion.
\end{quote}

\subsection{Gate attack summary}

\begin{enumerate}
\item \textbf{Gate 1: compact finite-first holomorphic \(E_3\) source.}
The manuscript uses \(\mathcal A_X^{E_3}\) and
\(\Phi^{\mathrm{FA}}_3(X)\) as supplied Stage-1 data, not as constructed
objects.  Vol. III gives a two-stage template
\(\Phi_d=\Spch\circ\Phi_d^{\mathrm{FA}}\), but its compact \(d=3\)
application to \(K3\times E\) still requires framing, descent, anomaly,
completion, and comparison witnesses.  The Hall atlas below cannot
manufacture the \(E_3\) source; it can only be compatible with it once
the source exists.

\item \textbf{Gate 2: Ran(\(E\)) elliptic-degree obstruction.}
Ordinary support locality on \(\Ran(E)\) sees only projection-finite
objects.  If a curve has positive elliptic degree, its support
dominates \(E\).  Therefore the atlas must be hybrid: local
projection-finite legs over disjoint opens of \(E\), global wrapped legs
for positive elliptic degree, and mixed extension stacks retaining
relative \(E\)-position before reduction.

\item \textbf{Gate 3: oriented compact Hall/CoHA atlas.}
This gate is open.  The needed atlas is not the OP scalar moduli space
and not the formal current envelope of \(\gDelta\).  It is an
extension-closed, HN-completed, oriented critical Hall object on compact
\(K3\times E\) sectors before \(E\)-translation reduction.  Its exact
data are listed in the next subsection.

\item \textbf{Gate 4: protected integration and orientation character.}
General critical CoHA technology supplies the language of vanishing
cycles, orientations, and motivic integration.  The compact protected
integration map for this \(K3\times E\) atlas is not constructed.  The
OP scalar theorem supplies an Euler-characteristic shadow only.

\item \textbf{Gate 5: \(\epsor=\nuDelta\).}
The automorphic character \(\nuDelta\) is theorem-level.  The Hall-side
character \(\epsor\) has no domain until the oriented completed Hall
object and a lifted type-II Weyl action are constructed.  The equality
is an orientation-monodromy theorem still to be proved.

\item \textbf{Gate 6: primitive recognition.}
The determinant fixes signed root supermultiplicities.  It does not fix
the compact Hall bracket, simple primitive representatives, full parity
dimensions, Borcherds--Kac relations, Hopf pairing, or radical quotient.
Primitive recognition remains conditional on presentation-level data.
\end{enumerate}

\subsection{Exact atlas data required}

Let \(X=S\times E\), where \(S\) is a \(K3\) surface and
\(p_E:X\to E\) is projection.  An oriented compact Hall/CoHA atlas
adequate for the Igusa realization must supply the following typed
objects.

\begin{enumerate}
\item \textbf{Charge and chamber data.}
A numerical charge lattice
\[
   \Gamma_X=K_0(\Perf(X))/\!\sim
\]
with Euler form, an Igusa degree map
\[
   \deg_{\Delta}:\Gamma_X\longrightarrow \La^{2,1}_{II},
\]
and a chamber map whose active coordinates are the Borcherds variables
\((n,l,m)\).  The map must separate the K3 norm, the elliptic degree,
and the \(D0\)-direction, and must send the protected root class to
\[
   \alpha(n,l,m)=2n f_2-lf_3+2m f_{-2}
\]
in the \(N=1\) normalization.

\item \textbf{Stability and finite-first control.}
A stability/HN datum
\[
   \sigma=(Z,\mathcal P),\qquad Q,\qquad S\subset\C^\times,
\]
with support form \(Q\), active-ray-free strict sector \(S\), and
finite saturated Hall-lower sets \(L_{\leq N,\leq R}\).  The finite
sets must be closed under HN factors, not merely under visible active
roots.  The finite quotient is only meaningful after these lower sets
are supplied.

\item \textbf{Moduli stacks before reduction.}
For every \(\gamma\in\Gamma_X\), a derived or d-critical stack
\[
   \mathfrak M^\sigma_\gamma(X)
      \subset \RPerf(X)_\gamma
\]
of semistable compact objects in the chosen sector.  This must be
defined before quotienting by translations of \(E\).  The stack must be
locally of finite type on each \(L_{\leq N,\leq R}\) and carry the
reduced obstruction theory or \((-1)\)-shifted critical structure
appropriate to the protected theory.

\item \textbf{Projection-finite and wrapped substacks.}
For opens \(U\subset E\),
\[
   \mathfrak M^{\mathrm{loc}}_\alpha(U)
   =
   \{A\in\mathfrak M^\sigma_\alpha(X):
     p_E(\Supp A)\subset U,\;p_E(\Supp A)\hbox{ finite}\},
\]
and wrapped substacks
\[
   \mathfrak M^{\mathrm{wr}}_\eta
   =
   \{A\in\mathfrak M^\sigma_\eta(X):
     (p_E)_*[\Supp_1 A]=b[E],\ b>0\}.
\]
The output of any extension involving a wrapped leg is wrapped.

\item \textbf{Critical charts and coefficients.}
For each finite stage, an oriented critical atlas
\[
   (\mathfrak U_i,f_i,G_i)_{i\in I}
\]
on the relevant stacks, with vanishing-cycle coefficients
\[
   \phi_{f_i}\otimes\mathscr L_{o_i}
\]
and transition isomorphisms on overlaps.  In shifted-symplectic
language this is the d-critical local presentation induced from the
PTVV \((-1)\)-shifted symplectic structure on the CY3 moduli stack,
together with the extra reduced-quotient correction.

\item \textbf{Extension correspondences.}
For local, wrapped, and mixed sectors, extension stacks of exact
triangles:
\[
\mathfrak M_{\gamma_1}^\sigma\times\mathfrak M_{\gamma_2}^\sigma
\xleftarrow{\ p_{\gamma_1,\gamma_2}\ }
\mathfrak E_{\gamma_1,\gamma_2}^\sigma
\xrightarrow{\ q_{\gamma_1,\gamma_2}\ }
\mathfrak M_{\gamma_1+\gamma_2}^\sigma.
\]
For disjoint opens \(U_i\subset E\), the local-local correspondence
must restrict to projection-finite supports.  The mixed correspondence
must have the form
\[
\prod_i\mathfrak M^{\mathrm{loc}}_{\alpha_i}(U_i)
\times\mathfrak M^{\mathrm{wr}}_\eta
\xleftarrow{\ p\ }
\mathfrak E^{\mathrm{hyb}}_{\alpha_1,\ldots,\alpha_r;\eta}
\xrightarrow{\ q\ }
\mathfrak M^{\mathrm{wr}}_{\eta+\sum_i\alpha_i}.
\]
The wrapped-wrapped correspondence is a global extension stack
\[
\mathfrak M^{\mathrm{wr}}_{\eta_1}\times
\mathfrak M^{\mathrm{wr}}_{\eta_2}
\xleftarrow{\ p\ }
\mathfrak E^{\mathrm{wr}}_{\eta_1,\eta_2}
\xrightarrow{\ q\ }
\mathfrak M^{\mathrm{wr}}_{\eta_1+\eta_2}.
\]

\item \textbf{Properness or Borel--Moore pull--push.}
The Hall product must be defined by a legitimate correspondence
formalism:
\[
   x*y
   =
   q_!\Bigl(\TS_o\bigl(p^*(x\boxtimes y)\bigr)\Bigr).
\]
Thus either \(q\) is proper on the finite HN stage, or one works in a
Borel--Moore/critical CoHA theory with compact-support conventions
which make the pushforward meaningful.  This is a theorem to prove for
the chosen compact sector, not a formal consequence of compactness of
\(X\).

\item \textbf{Strong orientation.}
For every object \(A\in\mathfrak M_\gamma^\sigma(X)\), a square root
of the virtual determinant line
\[
   \ell_A^{\otimes 2}
   \simeq
   \sdet \RHom(A,A),
\]
or equivalently a square root of the d-critical canonical line.  These
square roots must be compatible with direct sums, exact triangles,
overlaps of critical charts, reduced obstruction theory, and all triple
extension flags.  The transport must produce Thom--Sebastiani
isomorphisms coherent under both parenthesizations.

\item \textbf{Flag associativity.}
For triples of charges, flag stacks
\[
   \mathfrak{Flag}_{\gamma_1,\gamma_2,\gamma_3}^{\sigma}
\]
and their local/wrapped variants must identify
\((x*y)*z\) with \(x*(y*z)\) after vanishing cycles, orientation
transport, shifts, Tate twists, and completion.  The two composites
must agree on the nose in the finite quotient or by a specified
chain homotopy compatible with inverse limits.

\item \textbf{Equivariant \(E\)-descent.}
The translation action of \(E\) must be kept until after extension
correspondences are formed.  One then needs descent of vanishing-cycle
coefficients, orientation local systems, Hall products, and compact
support pushforwards through the reduced \(E\)-quotient.  Quotienting
first loses relative \(E\)-position and destroys the mixed local/global
Hall operation.

\item \textbf{HN completion.}
The finite stages assemble to a pro-object
\[
   \widehat{\Hall}_{X,\sigma,S}^{\mathrm{or}}
   =
   \varprojlim_{N,R}
   \Hall^{\mathrm{or}}_{X,\sigma,S,\leq N,\leq R},
\]
with continuous product and coproduct.  The inverse system must satisfy
Mittag--Leffler or an equivalent derived-limit control for products,
bar constructions, primitive projection, and protected integration.

\item \textbf{Comparison ports.}
The atlas must expose ports to the other gates:
\[
   \mathcal A_X^{E_3}\longrightarrow
   \mathcal F_X\longleftarrow
   \widehat{\Hall}_{X,\sigma,S}^{\mathrm{or}}
   \longrightarrow
   C_X,\quad
   \Prim_{\mathrm{prot}},\quad
   I^{\mathrm{prot}},\quad
   \rho_o.
\]
These ports are not optional decoration.  Without them the atlas may be
a Hall algebra, but it is not the compact Igusa realization datum.
\end{enumerate}

\subsection{First-principles construction route}

\begin{enumerate}
\item Start with the 3CY category \(\Perf(X)\), not with the scalar
partition function.  PTVV supplies the \((-1)\)-shifted symplectic
structure on the derived moduli of perfect complexes once the CY3
structure is fixed.  This gives the formal source of critical charts.

\item Choose the actual compact BPS sector.  The stable-pair sector,
ideal-sheaf sector, and full \(\Perf(X)\) sector are different.  The
sector must be extension-closed if it is to support a Hall product.
Stable pairs alone are usually not extension-closed in the needed
sense; they may be modules over a Hall algebra rather than the algebra
itself.

\item Fix stability and charge completion.  A finite active Borcherds
support is not an additive heart.  The finite object must be an HN
quotient by saturated Hall-lower sets.  This is exactly the
finite-first discipline in the manuscript and in Vol. III.

\item Build derived moduli stacks \(\mathfrak M_\gamma^\sigma\) and
extension stacks \(\mathfrak E_{\gamma,\eta}^\sigma\).  Verify finite
type in each quotient, the support property, and the proper or
Borel--Moore conditions needed for \(q_!p^*\).

\item Put the critical atlas and vanishing cycles on each finite stage.
The local models may be presented as critical loci of trace functions
only after the d-critical charts and their overlap maps have been
specified.  A slogan ``CY3 moduli are critical'' is not enough for a
global CoHA product.

\item Choose strong orientation data.  The Hall-side orientation torsor
may trivialize on favorable Kummer/DWR covers, but the relative
orientation class for hCS--Hall comparison and the Weyl monodromy sign
remain separate data.

\item Construct local, wrapped, and mixed extension correspondences.
Projection-finite objects give the ordinary \(\Ran(E)\) local part.
Positive elliptic degree gives global wrapped legs.  The mixed stacks
are the first genuinely compact innovation required by the Igusa
problem.

\item Prove associativity by flag stacks.  The proof must include the
two Thom--Sebastiani parenthesizations and the orientation transports.

\item Perform \(E\)-reduction only after Hall multiplication is defined.
Then prove descent of critical coefficients and orientations through the
reduced quotient.

\item Assemble the inverse limit and define the protected integration
map.  Its scalar shadow must recover the OP normalization, while its
primitive shadow must feed the Borcherds recognition theorem.
\end{enumerate}

\subsection{Obstruction ledger}

\begin{center}
\small
\begin{tabular}{p{0.27\textwidth}p{0.24\textwidth}p{0.39\textwidth}}
\hline
Obligation & Present status & Failure mode if omitted \\
\hline
Extension-closed compact sector &
Open &
No Hall product; stable-pair scalar spaces give modules or traces, not
the algebra. \\
Bridgeland/Hall stability with support property &
Open for full \(K3\times E\) Igusa sector &
No finite HN tower; infinite sums are not coefficientwise defined. \\
Finite-type semistable stacks in HN bounds &
Open &
Vanishing-cycle BM homology and motivic integration need finite-stage
control. \\
Proper/BM Hall correspondences &
Open &
The convolution \(q_!p^*\) may not exist or may not preserve the chosen
completion. \\
Reduced \(E\)-quotient descent &
Open &
Quotienting too early loses relative \(E\)-position and kills mixed
wrapped operations. \\
Hybrid local/wrapped extension stacks &
Open &
Ordinary \(\Ran(E)\) locality misses the \(s\)-direction. \\
Strong orientation with triple-extension coherence &
Partly available as general technology; not constructed here &
Products have undefined signs; no \(\epsor\). \\
hCS--Hall comparison primitives &
Finite Vol. III criterion only after supplied data &
No bridge from \(E_3\) observables to the Hall atlas. \\
Protected integration &
Open &
Scalar OP invariants do not become a Hall character or primitive trace. \\
Weyl lift &
Open &
\(\epsor(w)\) has no domain; comparison with \(\nuDelta\) is meaningless. \\
Primitive presentation certificate &
Open &
Same signed determinant can have the wrong bracket, parity, or radical. \\
\hline
\end{tabular}
\end{center}

\subsection{Healed conditional theorem}

\begin{theorem}[conditional oriented compact Hall/CoHA atlas]
Let \(X=K3\times E\).  Assume the following data are supplied.
\begin{enumerate}
\item An extension-closed compact sector of \(\Perf(X)\) with numerical
charge lattice \(\Gamma_X\), Igusa degree map to \(\La^{2,1}_{II}\),
and a stability/support/HN datum \((\sigma,Q,S)\).
\item For every finite saturated Hall-lower set \(L_{\leq N,\leq R}\),
finite-type semistable stacks
\(\mathfrak M^\sigma_{\gamma,\leq N,\leq R}\), their HN hulls, and
extension stacks of exact triangles.
\item A \((-1)\)-shifted critical or reduced perfect-obstruction atlas
on all local, wrapped, and mixed stacks, with vanishing-cycle
coefficients in a fixed motivic or Borel--Moore target.
\item Properness, or a stated Borel--Moore compact-support formalism,
which makes all Hall pull--push products on finite stages well-defined.
\item A strong orientation system
\(\ell_\gamma^{\otimes 2}\simeq\sdet\RHom(-,-)\), compatible with
overlaps, direct sums, exact triangles, Thom--Sebastiani, and triple
flags.
\item Projection-finite local substacks, wrapped positive
elliptic-degree substacks, and mixed local/wrapped extension
correspondences retaining relative \(E\)-position before reduction.
\item Equivariant descent through the reduced \(E\)-quotient after the
Hall product has been formed.
\item Compatibility of all products, orientations, and transition maps
under the inverse system \((N,R)\).
\end{enumerate}
Then the finite quotients form an associative inverse system of
oriented critical Hall/CoHA algebras
\[
   \Hall^{\mathrm{or}}_{X,\sigma,S,\leq N,\leq R},
   \qquad
   \widehat{\Hall}^{\mathrm{or}}_{X,\sigma,S}
   =
   \varprojlim_{N,R}
   \Hall^{\mathrm{or}}_{X,\sigma,S,\leq N,\leq R}.
\]
The projection-finite part gives the ordinary \(\Ran(E)\) local Hall
object.  The wrapped and mixed correspondences give the additional
global elliptic-degree operations required by the \(s\)-direction of
the Igusa product.  The resulting pro-object is the gate-3 atlas
required by the compact Igusa datum.
\end{theorem}

\begin{proof}[Proof sketch]
On each finite stage the Hall product is the vanishing-cycle
pull--push along the supplied extension correspondence, with the
Thom--Sebastiani and orientation transports inserted.  Associativity is
the two-Segal/flag-stack identity for length-three filtrations,
together with the assumed coherence of the orientation line.  The
saturated Hall-lower condition makes the complement a two-sided Hall
ideal, so finite quotients are closed under convolution.  The support
property and finite-type hypotheses make the finite stages legitimate
objects in the chosen coefficient theory.  The transition maps are
quotient maps by closed Hall ideals; hence the inverse limit carries a
continuous product.  Local/wrapped/mixed correspondences are separate
summands of the same extension-stack formalism, and the descent
hypothesis ensures that the reduced \(E\)-quotient is taken only after
relative position has served the Hall operation.

The theorem constructs no input datum.  It records the assembly theorem
which becomes available after the moduli stacks, orientations,
properness, wrapped correspondences, and descent maps have been
constructed.
\end{proof}

\subsection{Compatibility with the other gates}

\paragraph{Gate 1.}
The atlas must be a target of the compact \(E_3\) source, not a
substitute for it.  The comparison port is a map from the supplied
\(\mathcal A_X^{E_3}\) or hCS/BV observables to the oriented Hall
cosheaf.  Vol. III's DWR/Ran theorem proves sufficiency once a full
oriented hCS--Hall comparison datum is supplied; it does not supply the
compact \(E_3\) source.

\paragraph{Gate 2.}
The atlas must implement the hybrid repair.  Projection-finite
substacks support ordinary \(\Ran(E)\) locality.  Wrapped substacks and
mixed correspondences carry positive elliptic degree.  The report
rejects the weaker move ``ordinary Ran plus scalar Fock factor'' unless
the mixed local/wrapped Hall action and primitive brackets are also
constructed.

\paragraph{Gate 4.}
Protected integration is a further map out of the atlas:
\[
   I^{\mathrm{prot}}_{\sigma,o}:
   \widehat{\Hall}^{\mathrm{or}}_{X,\sigma,S}
   \longrightarrow
   \widehat{\C[\Gamma_{\mathrm{eff}}]}.
\]
It must be continuous, multiplicative, wall-crossing compatible, and
compatible with primitive projection.  Its scalar shadow may match OP;
that shadow does not define the map.

\paragraph{Gate 5.}
The orientation character requires both the strong orientation and a
lifted Weyl action
\[
   \rho_o:W^{(2)}(\La^{2,1}_{II})
   \to
   \operatorname{Aut}
   (\widehat{\Hall}^{\mathrm{or}}_{X,\sigma,S}).
\]
Only then is
\[
   \epsor(w)=\operatorname{sgn}_o(\rho_o(w))
\]
defined.  The desired equality
\[
   \epsor|_{W^{(2)}(\La^{2,1}_{II})}
   =
   \nuDelta|_{W^{(2)}(\La^{2,1}_{II})}
\]
is a determinant-line monodromy theorem, not a normalization choice for
\(\Delta_5\).

\paragraph{Gate 6.}
The primitive functor is applied to the completed Hall/factorization
object only after coproduct, augmentation, protected integration, parity
grading, and radical quotient are fixed.  The atlas must supply the
bracket constants; Gritsenko--Nikulin supplies the target presentation
and the denominator identity.  Equality of signed dimensions is a test,
not the proof.

\subsection{Local and cross-volume anchors}

\begin{itemize}
\item \verb|/Users/raeez/igusa-cusp-form/main.tex:903--933|:
protected Hall integration is stated as a criterion and explicitly not
constructed by general CoHA technology.
\item \verb|/Users/raeez/igusa-cusp-form/main.tex:1911--1957|:
the \(E_3\) upgrade begins with a supplied Stage-1 object and supplied
\(K3\)-to-\(E\) specialization.
\item \verb|/Users/raeez/igusa-cusp-form/main.tex:2018--2087|:
projection-finite versus wrapped elliptic-degree sectors and the two
repair paths.
\item \verb|/Users/raeez/igusa-cusp-form/main.tex:2106--2181|:
sectorial Hall truncation and conditional finite Hall quotient.
\item \verb|/Users/raeez/igusa-cusp-form/main.tex:2183--2237|:
compact Igusa realization datum, including the oriented Hall object.
\item \verb|/Users/raeez/igusa-cusp-form/main.tex:2265--2361|:
primitive recognition is conditional on the compact
Hall/factorization construction.
\item \verb|/Users/raeez/igusa-cusp-form/main.tex:2408--2453|:
the determinant does not determine Hall constants, primitive
representatives, parity, or relations.
\item \verb|/Users/raeez/igusa-cusp-form/proj.bib:230--276|:
Thomas, Behrend, Kontsevich--Soibelman, and
Davison--Meinhardt are present as the DT/critical-CoHA technology base.
\item \verb|agent_material/14_compact_E3_source_from_volumes_attack_heal.tex|:
the compact \(E_3\) source is not present in the surrounding volumes.
\item \verb|agent_material/14_ranE_hybrid_wrapping_from_volumes_attack_heal.tex|:
ordinary \(\Ran(E)\) locality misses positive elliptic degree; hybrid
local/global wrapping is the correct target.
\item \verb|agent_material/14_orientations_protected_integration_from_volumes_attack_heal.tex|:
\(\epsor\) is undefined without the oriented Hall object and Weyl lift.
\item \verb|agent_material/14_primitive_borcherds_recognition_from_volumes_attack_heal.tex|:
primitive recognition requires presentation-level data, not signed
Euler characteristics.
\item \verb|agent_material/14_chiral_bar_cobar_Koszul_realization_attack_heal.tex|:
bar--cobar inversion is downstream of a supplied curve-level object.
\item \verb|/Users/raeez/calabi-yau-quantum-groups/chapters/theory/bps_positive_geometry_closure.tex:1--170|:
Vol. III finite-first oriented motivic Hall cosheaf formalism; it
requires chamber data, support, HN control, orientation, and admissible
correspondences.
\item \verb|/Users/raeez/calabi-yau-quantum-groups/chapters/theory/bps_positive_geometry_closure.tex:231--330|:
the compact chamber certificate and Igusa quotient boundary; the raw
Hall theorem is stronger than the automorphic quotient.
\item \verb|/Users/raeez/calabi-yau-quantum-groups/chapters/theory/bps_positive_geometry_closure.tex:334--385|:
seven-class hCS--Hall descent obstruction, including orientation,
Thom--Sebastiani, factorization, compact support, and completion.
\item \verb|/Users/raeez/calabi-yau-quantum-groups/chapters/theory/cy3_chain_level_bridge.tex:522--700|:
oriented hCS--Hall comparison datum on the DWR/Ran nerve; the theorem
starts after the full datum is supplied.
\item \verb|/Users/raeez/calabi-yau-quantum-groups/chapters/theory/cy3_chain_level_bridge.tex:1099--1211|:
descent obstruction for the oriented comparison map.
\item \verb|/Users/raeez/calabi-yau-quantum-groups/chapters/theory/cy3_chain_level_bridge.tex:3452--3705|:
Hall orientation torsor and finite DWR/Ran comparison for \(K3\times E\);
global pro-compatibility and relative comparison data remain required.
\item \verb|/Users/raeez/calabi-yau-quantum-groups/chapters/examples/k3e_bkm_chapter.tex:13977--14136|:
Vol. III separates the theorem-grade character
\(-\Delta_5^{-2}\) from the conditional compact critical-CoHA reading.
\item \verb|/Users/raeez/chiral-bar-cobar/worldview_synthesis_2026_04_17.tex:740--754|:
the six routes to \(G(K3\times E)\) do not produce a common algebra
without the missing monoidal equivalence and Hall package.
\item \verb|/Users/raeez/chiral-bar-cobar-vol2/main.tex:1512--1540|:
Vol. II records the scalar \(K3\times E\) shadow and the missing
gravity-line Hall--Borcherds comparison.
\item \verb|/Users/raeez/topological-strings/frontier_mnop_framing_volume.tex:125--145|:
Joyce/KPS orientation technology gives generic compact-CY3 vertex and
orientation inputs; it does not construct the Igusa Hall atlas.
\end{itemize}

\subsection{Status recommendation}

\begin{enumerate}
\item The paper should not say that the compact Hall/CoHA atlas is
constructed.
\item The compact section should present gate 3 as the central
construction problem with the exact atlas data above.
\item The strongest true theorem available now is the conditional atlas
assembly theorem: supplied moduli, critical charts, orientation,
properness, wrapped correspondences, and descent imply a finite-first
oriented pro-Hall object.
\item The scalar OP identity and the Borcherds denominator should remain
the theorem core.  The compact Hall atlas is a frontier programme
attached to that core.
\item The next actual mathematical task is not more exposition.  It is
to choose the compact sector and prove the first finite HN atlas:
charges, stability, finite-type stacks, extension correspondences, and
orientation on one cofinal tower.
\end{enumerate}

\subsection{Files changed}

Only this file was changed:
\[
\texttt{agent\_material/15\_oriented\_compact\_Hall\_CoHA\_atlas\_attack\_heal.tex}.
\]

\subsection{Checks run}

\begin{itemize}
\item Read \verb|CLAUDE.md| and \verb|AGENTS.md|.
\item Read the relevant compact-realization, Hall, orientation,
Ran(\(E\)), and primitive-recognition anchors in \verb|main.tex|.
\item Read \verb|proj.bib| anchors for DT, Behrend, KS CoHA,
Davison--Meinhardt, and Costello--Gwilliam technology.
\item Read all six \verb|agent_material/14_*_attack_heal.tex| reports.
\item Read targeted cross-volume anchors in
\verb|~/calabi-yau-quantum-groups|,
\verb|~/chiral-bar-cobar|,
\verb|~/chiral-bar-cobar-vol2|, and
\verb|~/topological-strings|.
\item Verified before writing that the target file did not exist.
\end{itemize}

\subsection{Remaining open questions}

\begin{enumerate}
\item Which exact compact sector is intended: full \(\Perf(K3\times E)\),
stable sheaves, stable pairs, an \(E\)-reduced quotient category, or a
hybrid Hall category acting on stable-pair modules?
\item Does the needed Bridgeland or Hall stability condition exist on
the full Igusa charge sector with a support property strong enough for
the HN finite tower?
\item Are the semistable stacks in each finite HN lower set finite type,
and is the extension correspondence proper or admissible in the chosen
Borel--Moore theory?
\item Can the reduced \(E\)-quotient be postponed until after Hall
multiplication without breaking the OP scalar comparison?
\item What is the first nontrivial mixed local/wrapped extension stack,
and what is its orientation line?
\item Does the Hall-side orientation torsor trivialization on the
Kummer/DWR cover descend to the actual compact sector and survive the
pro-limit?
\item Can the type-II Weyl reflections be lifted to autoequivalences or
wall-crossing functors of the oriented completed Hall object?
\item At the first finite height, can one exhibit primitive classes for
the three real simple roots and compute their Hall brackets?
\end{enumerate}
